\Problem{Els campanars de Sant Esteve Sesrovires}

\Statement

\FigureR{width=4.5cm}{campanar}

L'església de Sant Esteve Sesrovires té la peculiaritat de tenir dos campanars que coronen l'església.

\medskip 

La llegenda explica que fa molts anys hi havia un matrimoni molt ric de la població que desitjaven molt tenir fills, però passaven els anys i no aconseguien fer realitat els seus desitjos. En aquells temps la població estava immersa en construir l'es\-glé\-sia nova, feta tota amb totxanes de terrissa. Els vilatans coïen les totxanes en un lloc anomenat la Vinya del Forn, que era on hi havia un gran forn per coure. Feien cadena tots i de mà en mà portaven les totxanes a la plaça on s'estava construint l'església.

\medskip 

Durant la construcció, el matrimoni que també ajudava en la tasca, va demanar a Déu que si els donava un fill, tant era que fos hereu o pubilla, ells pagarien la construcció d'un campanar alt i majestuós per damunt de les cases de la població i que es pogués veure de les poblacions veïnes. Al cap de poc temps la dona es va adonar que estava embarassada i quan va arribar l'hora del naixement, no van tenir un fill sinó que en van tenir dos. I en senyal de gratitud van complir la seva promesa, però en comptes de pagar un sol campanar, en van construir dos per agraïr el naixement dels bessons.

\medskip

En temps més moderns, es va decidir incloure un rellotge a cadacun dels campanars. Feu un programa que pinti els dos campanars a l'hora indicada, de forma semblant al problema P48022.

\Sample